% Topic-keyword dyads used to filter documents by issue area.
% Rebuilt from the table in "Thesis Abridged.docx" (Figure 6 there; Figure 11 in the full thesis).
\begin{table}[htbp]
\centering
\begin{threeparttable}
\caption{Topic--keyword dyads in US--China discourse}
\label{tab:topic-keywords}
\small
\begin{tabular}{@{}l l p{3.4in}@{}}
\toprule
 & Issue area & Keywords \\
\midrule
1 & Trade and finance & trade, tariff, protectionist, WTO, import, export, countervailing duty, intellectual property, manufactur(e/er/ing), jobs, exchange rate, currency (manipulation) \\
2 & Chinese domestic periphery & human rights, Xinjiang, Uyghur, Tibet, Hong Kong \\
3 & Taiwan & Taiwan \\
4 & Multilateral negotiations & denuclearize(ation), DPRK, North Korea, Iran, JCPOA, IAEA, Paris, climate, COP, ASEAN, APEC, EAS, G20, summit, multilateral \\
5 & South China Sea & South China Sea \\
6 & Global health and epidemics & COVID, coronavirus, pandemic, epidemic, outbreak, Ebola, SARS \\
7 & Russia & Russia \\
\bottomrule
\end{tabular}
\begin{tablenotes}[flushleft]
\footnotesize
\item \textit{Notes:} A document is assigned to every topic whose keywords it contains, so a document can belong to more than one topic.
\end{tablenotes}
\end{threeparttable}
\end{table}
